\clearpage
\fancyhead[L]{\small BODE. ASTRONOMICAL YEARBOOK.}
\hypertarget{native-568-lower}{}
\begin{center}\small Göttingische gelehrte Anzeigen. No. 40. Pages 393..400. 10 March 1814.\\[3pt]\rule{.72\linewidth}{.3pt}\end{center}

\emph{Astronomical Yearbook for the Year} 1816, \emph{together with a collection of the latest papers, observations and notices relating to the astronomical sciences. Calculated and edited by} J.\,E. \textsc{Bode}, Royal Astronomer and member of the Academy. 1813. 268 pages in octavo.

The year 1816 is marked by a great solar eclipse (19 November), which will have a period of totality in part of Europe, for example in Danzig and Warsaw. The calendar is arranged in exactly the same way this year as in the preceding one. The latest determinations now serve as the basis for the Sun's diameter and the obliquity of the ecliptic. For the eclipses of Jupiter's satellites, which are still calculated from \textsc{Wargentin}'s tables, a general account of how the differences between those tables and \textsc{Delambre}'s vary has been added. We were all the more pleasantly surprised by the number and richness of the \emph{supplements} that accompany this volume too, since in the circumstances of the time we had scarcely hoped to see the Yearbook appear at the usual time at all. It opens with a chronological list of the most celebrated (deceased) astronomers since the thirteenth century, with their achievements, writings and discoveries. One would be hard pressed to find a name of any note in the history of astronomy missing from it. — Astronomical observations made at the Royal Observatory in Copenhagen in 1811 and 1812 by Mr. \textsc{Bugge}, Councillor of State and Knight. Oppositions of Uranus in 1811 and 1812, of Saturn in 1811 and 1812, of Ceres in 1811, of Vesta in 1812 and of Mars in 1811, observations of the great comet of 1811, stellar occultations and eclipses of Jupiter's satellites. — Eclipses of Jupiter's satellites and stellar occultations observed at Greenwich Observatory in 1809 and 1810 (taken from the printed collection of \textsc{Maskelyne}'s observations). — Observations of the variable star $\eta$ in Antinous and tables for calculating its maximum brightness, by Prof. \textsc{Wurm}, a continuation of the paper in the Yearbook for 1814. We wish this distinguished astronomer would soon give us a similar work on the variable star in Cetus. — Astronomical observations made at the Royal Observatory in Berlin in 1812. — Observations of Jupiter's satellites. — Eclipses, stellar occultations, oppositions of Ceres, Pallas, Juno, Vesta and the comet of 1812, observed at Vienna Observatory by Mr. \textsc{Triesnecker}. Meridian observations of Ceres and Vesta, of Mars, Saturn and Jupiter, made in 1811 by Mr. \textsc{Groombridge}. — On determining the graduation errors of a reflecting sextant, by Dr. \textsc{Benzenberg}. There is good reason for saying that the absolute magnitudes of angles measured with these instruments cannot be relied on without qualification, even if in \SourceLeaf{569}individual cases it may remain undecided how much of this uncertainty is due to actual graduation errors, to a slight eccentricity, or to a slight departure from parallelism of the two faces of the large mirror. The method Mr. B. proposes for tabulating the errors of the instrument can, however, be used only in very level country unless one also has an instrument for measuring small altitudes. It may also be worth urging special care to ensure that all the angles are measured from exactly the same point or at least reduced to that point. — Various astronomical observations communicated by Mr. \textsc{Lee} from London. Observations of the great comet of 1811 in the East Indian Ocean aboard a ship from 17 May to 15 June, and various astronomical observations from Port Jackson in New Holland and from Calcutta. — Astronomical observations made in 1812 at Prague Observatory by Messrs. \textsc{David} and \textsc{Bittner}. — Observed zenith distances of the Sun and stellar occultations at St. Gallen in Switzerland, by Mr. \textsc{von Scherer}. Mr. \textsc{Triesnecker} also derives the geographical position of the place from these observations. — Visible changes in the brightness of Algol in 1814, 1815 and 1816, predicted by Mr. \textsc{Wurm}. — Observations of the great comet of 1811 from 6 September to 19 October, oppositions of Saturn in 1811 and stellar occultations by Mr. \textsc{Derflinger} at Kremsmünster. — New table of refraction derived from \textsc{Bradley}'s observations of circumpolar stars, by Prof. \textsc{Bessel}. The table is taken from a paper by the author in the fourth issue of the Königsberg Archive. Instructions for its use are missing, and are all the more needed because the worked example does not quite follow the author's approach. — Observations of the great comet in 1811, with remarks on the structure of its various parts, by Dr. \textsc{Herschel}. Later come observations of the second comet of 1811, with remarks on its structure, by the same author. These are extracts from two papers in the latest volumes of the \emph{Philosophical Transactions}. We can all the more readily dispense with a detailed account of them here because the original papers have already been reviewed in these pages by another reviewer. We therefore remark only that, valuable as \textsc{Herschel}'s observations of these comets are, it still seems to us much too early to try to establish conclusions about the natural history of these remarkable celestial bodies. There are still far too few observations, and even these are far less reliable than astronomical observations of other kinds. We are convinced that only future centuries, which will see the comets now observed return, will be able to decide what can now be only a vague hypothesis. Among other things, Mr. \textsc{Herschel} conjectures that the comet of 1807, before reaching perihelion on this occasion, had once approached another fixed star and so had, as it were, matured earlier than the comet of 1811. He believes that the possibility of comets orbiting suns other than our own is made very probable by the fact that, of the large number observed so far, we know with certainty of the return of only one. In this, however, he evidently overlooks the fact that, according to all results to date, the orbital periods of comets generally have to be reckoned in millennia rather than in centuries, that comets have been properly observed for only a couple of centuries, and that their orbits have been calculated for little more than one century. As for the comet of 1807 in particular, we know that its greatest distance from the Sun, like the greatest distance of every other comet whose orbit could be calculated as an ellipse, is next to nothing compared with the distance of the nearest fixed stars. — Stellar occultations, eclipses of Jupiter's satellites, a lunar eclipse, Vesta, Uranus and Mars, observed at the observatory in Wilna in 1811, 1812 and 1813 by Mr. \textsc{Joh. Sniadecky}. — Observations of the great comet of 1811 at the observatory in Palermo and calcu\SourceLeaf{570}lation of the elements of its orbit by Mr. \textsc{Piazzi}, from his printed paper. — From a letter by Prof. \textsc{David} in Prague. It concerns the differing results that Messrs. \textsc{Vent}, \textsc{Bode} and \textsc{David} found at different times for the altitude of the pole at the Bohemian Riesenkuppe, on which we are of course in no position to pass judgment here. — Supplement by Mr. \textsc{von Wisniewsky} to the list in the Yearbook for 1813 of stellar occultations visible in 1813. — List of the longitudes and latitudes of nine of the principal fixed stars according to the latest observations, taken from the \emph{Nautical Almanac} for 1815. — Geographical position of various places in Central America, determined by Mr. \textsc{von Humboldt} (from his \emph{Recueil d'observations astronomiques etc.}). — Method of projection for a general map of the sky. It makes distances from the north pole in the projection proportional to the tangents of $\frac{2}{5}$ of the distances on the sphere, whereas in stereographic projection, as is well known, they are proportional to the tangents of half the distances on the sphere. The author used this new method to draw the sky map for the latest edition of his esteemed guide to the starry sky, earlier editions of which included a similar map made by stereographic projection. He chose this modification so that the degree intervals toward the south pole would not grow so disproportionately large compared with those in the north. Of course, this gives up several advantages of stereographic projection. The image is no longer similar to the original \emph{in its smallest parts}. All circles lying obliquely to the equator are represented by curves of higher order rather than by circles, and the curves can be constructed only with difficulty, point by point. The angles at which they intersect are no longer equal to the angles on the sphere. — Geocentric path of Pallas from 1 August 1814 to 1 February 1815, and of Vesta from 18 October 1813 to 9 July 1814. — Observations of planets at Greenwich Observatory in 1809 and 1810 by Dr. \textsc{Maskelyne}. — Observations of the two comets that appeared in 1813. Only four observations of the first comet from Paris are listed here. For the other, discovered by Prof. \textsc{Harding} at our observatory and at the same time by Mr. \textsc{Pons} in Marseille, only some general information is given. — Elliptical elements of the planetary orbits, those for the older planets from the third edition of \textsc{Laplace}'s \emph{exposition du système du monde} (a fourth edition appeared in 1812), and those for the newer planets from the latest determinations by \textsc{Gauss} known to the editor at the time. This article could already do with an almost complete revision, since the theories of most of the planets have been reworked in the meantime. — Further observations of the great comet of 1811, with astronomical remarks by Prof. \textsc{Bessel}, a new reduction of \textsc{von Zach}'s important observations of this comet on the first branch of its orbit. — Observations of the comet of 1812, with the elements of its orbit, by \textsc{Bouvard}. — Method for determining the deviation of a transit instrument from the meridian from observations of the upper and lower culminations of two northern stars nearly opposite in right ascension, with a list of such stars. We consider this the most suitable method of all, especially when one of the stars is Polaris itself, and remark only that in this case the other star may just as well have nearly the \emph{same} right ascension without any disadvantage. — Observations of planets at Paris Observatory in 1809 by Mr. \textsc{Bouvard} (from the \emph{Connoissance des tems} for 1812). — Of the various short astronomical notices added at the end, we single out here only those concerning the excellent astronomical instruments made in \textsc{Reichenbach}'s workshop in Munich, including a 24-foot achromatic telescope with an 8-inch aperture. — Two supplements to this volume are also worth noting. The first contains the original observations of the great comet of 1811 made by \textsc{Wisniewsky} in August 1812, the other \textsc{Pond}'s catalogue of fixed stars (though not \SourceLeaf{571}the very latest one). — We cannot help thanking the distinguished editor for the tireless zeal with which he has brought together so much material of importance to the astronomer. Even the material taken from other publications, most of which are now in any case inaccessible to the vast majority of readers, cannot but be welcome.
